%% ma: L12-B08-0093
%% lop: 12
%% chuong: 2
%% bai: 8
%% dang: 12.8.29
%% loai: TN4
%% muc: NB
%% dap_an: "B"
%% nguon: "CĐ2.3.pdf giáo viên gửi 27/10/2026 (M07, Lê Minh Khôi – Nguyễn Anh Khôi, Chương II Chủ đề 3 Biểu thức tọa độ của các phép toán vectơ), Câu 1"
%% kien_thuc: "Bài 8 (biểu thức tọa độ của phép cộng, trừ vectơ, tích vectơ với một số)"
%% trang_thai: chinh-thuc
\begin{ex}
Trong không gian $Oxyz$, cho hai vectơ $\vec a=(2;3;5)$ và $\vec b=(0;1;3)$. Tọa độ của vectơ $\vec u=\vec a+\vec b$ là
\choice{$\vec u=(-2;-2;-2)$}{\True $\vec u=(2;4;8)$}{$\vec u=(1;2;4)$}{$\vec u=(-2;-4;-8)$}
\loigiai{$\vec u=(2+0;3+1;5+3)=(2;4;8)$.}
\end{ex}
